\unitopen{08}
\begin{candos}
\item zamówić jedzenie i picie grzecznie (\de{Ich möchte \dots}, \de{Ich hätte gern \dots});
\item zapytać o cenę i zapłacić (\de{Was kostet das?}, \de{Zahlen, bitte!});
\item powiedzieć, co lubisz, co Ci smakuje i czego nie jesz;
\item czytać proste menu i rozumieć ceny w euro i frankach.
\end{candos}

\section{Słownictwo}
\vocabtable{08}

\section{\de{möchten}: zamawianie}
\begin{center}\small
\begin{tabular}{@{}ll@{}}
\de{ich möchte} & \de{wir möchten}\\
\de{du möchtest} & \de{ihr möchtet}\\
\de{er/sie/es möchte} & \de{sie/Sie möchten}\\
\end{tabular}
\end{center}
\de{möchte} nie ma końcówki w pierwszej i trzeciej osobie. Po \de{möchte} drugi czasownik jest na \emph{końcu} w
bezokoliczniku: \de{Ich möchte einen Kaffee trinken.} Najczęściej jednak opuszczasz drugi czasownik:
\de{Ich möchte einen Kaffee.} Alternatywa uprzejmiejsza: \de{Ich hätte gern \dots} (dosłownie „miałbym chętnie”).

\section{\de{essen}, \de{trinken}, \de{schmecken}}
\de{essen} jest nieregularny: \de{ich esse}, \de{du isst}, \de{er isst}, \de{wir essen}. \de{trinken} jest regularny.
\de{schmecken} łączy się z celownikiem i opisuje wrażenie: \de{Das Essen schmeckt gut.} \de{Die Suppe schmeckt
mir nicht.} (z \de{mir} za dwa tygodnie).

\section{Ceny}
\de{Was kostet das?} \de{Das kostet drei Euro fünfzig.} Zapis \de{3,50 €}. W Szwajcarii: \de{CHF 12.50} =
\de{zwölf Franken fünfzig}. Liczby powyżej 100: \de{hundertzwanzig}, \de{zweihundertfünfzig}, \de{tausend},
\de{eintausendzweihundert}. Rachunek: \de{Das macht zusammen 24 Euro.}

\begin{zasada}[Rodzajniki w zamówieniu]
W zamówieniu napoje i dania są dopełnieniem w bierniku: \de{Ich möchte \textbf{einen} Kaffee} (\de{der Kaffee}),
\de{eine} Suppe, \de{ein} Bier, \de{ein} Wasser. Dla kilku: \de{zwei Kaffee}, \de{drei Bier} (bez końcówki mnogiej po liczbie).
\end{zasada}

\begin{falle}
\de{Ich will einen Kaffee} brzmi jak rozkaz i bywa odbierane niegrzecznie. Polskie „chcę” zostaw w polskim:
używaj \de{möchte} albo \de{hätte gern}. Drugi błąd: \de{Bitte ein Kaffee}: po \emph{bitte} rodzajnik męski w bierniku,
czyli \de{einen Kaffee}. Trzeci: \de{Ich esse Kaffee} (Kawę się \emph{pije}: \de{trinken}).
\end{falle}

\begin{szwajcaria}
W Szwajcarii obsługa jest wliczona w cenę, ale napiwek (\de{Trinkgeld}) zaokrągla się zwykle w górę;
mówisz \de{Stimmt so} (reszty nie trzeba). Rachunek: \de{Zahlen, bitte} albo \de{Ich möchte bezahlen}.
\de{Das Brötchen} to w Zurychu \de{das Brötli}. Wodę z kranu poprosisz słowami \de{Leitungswasser} (może być
płatna w restauracji). Ceny w restauracjach są wyraźnie wyższe niż w Polsce; ucz się
porównywać w \emph{frankach}.
\end{szwajcaria}

\begin{praca}
Lunch z zespołem: \de{Was empfehlen Sie?} \de{Ich nehme das Tagesmenü.} \de{Wir zahlen getrennt.} /
\de{Wir zahlen zusammen.} \de{Ich lade dich ein} (zapraszam, płacę). Wybierając: \de{Ich bin Vegetarier} /
\de{Ich esse kein Fleisch.}
\end{praca}

\section{Dialogi}
\textbf{Dialog 1 --- w restauracji}
\begin{dialog}
\turn{A}{Guten Abend. Haben Sie einen Tisch für zwei?}{Dobry wieczór. Czy jest stolik dla dwóch osób?}
\turn{B}{Ja, gern. Hier ist die Speisekarte.}{Tak, proszę. Oto menu.}
\turn{A}{Ich hätte gern die Suppe und ein Wasser.}{Poproszę zupę i wodę.}
\turn{B}{Und was möchten Sie trinken?}{A co Państwo chcą pić?}
\turn{A}{Ich möchte einen Tee, bitte.}{Poproszę herbatę.}
\turn{B}{Möchten Sie noch etwas?}{Czy chcą Państwo coś jeszcze?}
\turn{A}{Nein, danke. Zahlen, bitte!}{Nie, dziękuję. Poproszę rachunek.}
\turn{B}{Das macht 18 Euro fünfzig.}{To razem osiemnaście euro pięćdziesiąt.}
\turn{A}{Stimmt so.}{Reszty nie trzeba.}
\end{dialog}

\section{Ćwiczenia}
\begin{cwiczenia}
\exitem{Uzupełnij: \de{Ich \dots{} einen Kaffee.} (möchten)}{\de{möchte}}
\exitem{Uzupełnij: \de{Du \dots{} gern Suppe.} (essen)}{\de{isst}}
\exitem{Uzupełnij: \de{Wir \dots{} zwei Bier.} (möchten)}{\de{möchten}}
\exitem{Uzupełnij rodzajnik: \de{Ich möchte \dots{} Kaffee.}}{\de{einen}}
\exitem{Uzupełnij rodzajnik: \de{Ich hätte gern \dots{} Suppe.}}{\de{eine}}
\exitem{Uzupełnij rodzajnik: \de{Ich möchte \dots{} Wasser.}}{\de{ein}}
\exitem{Którego czasownika użyjesz przy kawie: \de{essen} czy \de{trinken}?}{\de{trinken}}
\exitem{Zamów grzecznie herbatę, używając \de{hätte gern}.}{\de{Ich hätte gern einen Tee.}}
\exitem{Zapytaj o cenę.}{\de{Was kostet das?}}
\exitem{Przeczytaj cenę: 12,50 €.}{\de{zwölf Euro fünfzig}}
\exitem{Przeczytaj cenę: CHF 8.30.}{\de{acht Franken dreißig} (CH: \de{dreissig})}
\exitem{Poproś o rachunek.}{\de{Zahlen, bitte!}}
\exitem{Przetłumacz: „Jestem wegetarianinem, nie jem mięsa.”}{\de{Ich bin Vegetarier, ich esse kein Fleisch.}}
\end{cwiczenia}

\weekend{Odwiedź (albo zasymuluj z AI) restaurację: poproś o stolik, zamów danie i napój, zapytaj o cenę,
poproś o rachunek. Potem wykonaj rolę kelnera i odwróć dialog. Nagraj całość i sprawdź
rodzajniki w bierniku.}
